\documentclass{article}
\usepackage[utf8]{inputenc}
\usepackage{amsmath,amssymb}
\usepackage[most]{tcolorbox}
\usepackage{geometry}
\geometry{margin=2.5cm}

\newcommand{\step}[1]{\par\noindent\hangindent=3.5em\hangafter=1%
  \makebox[3.5em][l]{\textbf{#1.}}}
\newcommand{\steptag}[1]{\hfill{\small\textsf{[#1]}}}

\newtcolorbox{proofbox}[1]{
  colback=gray!5, colframe=gray!60,
  fonttitle=\bfseries, title={#1},
  breakable, enhanced,
  left=4pt, right=4pt, top=4pt, bottom=4pt,
  before skip=10pt, after skip=10pt
}

\begin{document}

\begin{proofbox}{Parameterized same-map quotient local-index and ambient-c...}
\step{1} Parameterized same-map quotient local-index and ambient-chart package: for every universal def-stage1-polar-witness-data tuple W=(C\_rect,C\_ch,C\_pol,C\_grp,C\_path,C\_der,e\_rect,kappa\_ch,kappa\_pol,kappa\_der,delta\_*,epsilon\_*\textasciicircum{}r,e\_S1,r\_iso) with C\_rect,C\_ch,C\_pol,C\_grp,C\_path,C\_der>=1, 0<e\_rect<=1/C\_rect, 0<kappa\_ch,kappa\_pol,kappa\_der<=1/2, delta\_*=min\{1/4,kappa\_ch/(4*C\_ch),kappa\_pol/(4*C\_pol)\}, epsilon\_*\textasciicircum{}r=min\{1/4,kappa\_ch/(4*C\_ch),kappa\_pol/(4*C\_pol),kappa\_der/(8*C\_der),1/C\_grp,delta\_*/(12*C\_path*C\_grp)\}, e\_S1=min\{e\_rect,epsilon\_*\textasciicircum{}r/C\_rect\}, and r\_iso=min\{delta\_*/4,kappa\_der/(8*C\_der)\}, there exist universal e\_quot\textasciicircum{}r>0 and epsilon\_B\textasciicircum{}r>0 with epsilon\_B\textasciicircum{}r=min\{epsilon\_*\textasciicircum{}r,e\_quot\textasciicircum{}r\} such that, for every finite-dimensional exact-unit epsilon\_r-C*-algebra with 0<=epsilon\_r<=epsilon\_B\textasciicircum{}r and 1<dim\_C calX<infinity, and for displayed data satisfying all of the following: (A\_2) for every V in calUbar\_delta\_* and A\textasciicircum{}par in B\_\{2delta\_*\}\textasciicircum{}\{icalH\}(0) there is a unique g\_V(A\textasciicircum{}par) in B\_\{2delta\_*\}\textasciicircum{}\{calH\}(0) with f\_V(A\textasciicircum{}par+g\_V(A\textasciicircum{}par))=0, where f\_V(A)=(1/2)*(((J+A\textasciicircum{}dagger) bold-dot V\textasciicircum{}dagger) bold-dot (V bold-dot (J+A))-J), chi\_V(A\textasciicircum{}par)=V bold-dot (J+A\textasciicircum{}par+g\_V(A\textasciicircum{}par)) lies in calU, ||Dg\_V(A\textasciicircum{}par)||<=C\_ch*(epsilon\_r+delta\_*), ||D\_\{A\textasciicircum{}perp\}f\_V(A\textasciicircum{}par+g\_V(A\textasciicircum{}par))-I\_calH||<=C\_ch*(epsilon\_r+delta\_*)<1, and these C\textasciicircum{}1 charts cover calU; (A\_4) set Pi\_delta\_*:calU x B\_\{delta\_*\}\textasciicircum{}\{calH\}(J)->calX by Pi\_delta\_*(U,K):=U bold-dot K and set S\_delta\_*:=Pi\_delta\_*(calU x B\_\{delta\_*\}\textasciicircum{}\{calH\}(J)); Pi\_delta\_* is a C\textasciicircum{}1 diffeomorphism onto S\_delta\_* with unique inverse maps u\_delta:S\_delta\_*->calU and h\_delta:S\_delta\_*->B\_\{delta\_*\}\textasciicircum{}\{calH\}(J) satisfying X=u\_delta(X) bold-dot h\_delta(X) for every X in S\_delta\_* and u\_delta(U)=U and h\_delta(U)=J for every U in calU; (A\_5) the displayed maps mu:calU x calU->calU and sigma:calU->calU are defined by mu(U,V):=u\_delta(U bold-dot V) and sigma(U):=u\_delta(U\textasciicircum{}dagger), are global C\textasciicircum{}1 and, for every U,V,Z in calU, mu(J,U)=mu(U,J)=U, sigma(J)=J, ||mu(U,V)-U bold-dot V||<=C\_grp*epsilon\_r, ||sigma(U)-U\textasciicircum{}dagger||<=C\_grp*epsilon\_r, ||mu(mu(U,V),Z)-mu(U,mu(V,Z))||<=C\_grp*epsilon\_r, ||mu(sigma(U),U)-J||<=C\_grp*epsilon\_r, and ||mu(U,sigma(U))-J||<=C\_grp*epsilon\_r; (A\_6) for every U\_0,U\_1 in calU and q in [0,1] satisfying ||U\_1-U\_0||<=q, C\_path*q<=1/4, and C\_path*(q+epsilon\_r*q+q\textasciicircum{}2)<delta\_*-C\_pol*(epsilon\_r*delta\_*+delta\_*\textasciicircum{}2), the path H(-,U\_0,U\_1):[0,1]->calU given by H(t,U\_0,U\_1):=u\_delta((1-t)*U\_0+t*U\_1) is defined, is jointly continuous in its displayed variables, and joins U\_0 to U\_1, and satisfies H(t,cU\_0,cU\_1)=c*H(t,U\_0,U\_1) for every c in U(1) and t in [0,1]; (A\_7) for every s in \{+1,-1\}, set chi\_s:B\_\{r\_iso\}\textasciicircum{}\{icalH\}(0)->calU by chi\_s(A):=sJ bold-dot (J+A+g\_\{sJ\}(A)), let phi\_\{sJ\}\textasciicircum{}par:chi\_s(B\_\{r\_iso\}\textasciicircum{}\{icalH\}(0))->B\_\{r\_iso\}\textasciicircum{}\{icalH\}(0) be its inverse, and set F\_s:B\_\{r\_iso\}\textasciicircum{}\{icalH\}(0)->icalH by F\_s(A):=phi\_\{sJ\}\textasciicircum{}par(sigma(chi\_s(A))); then sigma maps chi\_s(B\_\{r\_iso\}\textasciicircum{}\{icalH\}(0)) into itself and ||D(F\_s-id)(A)+2I||<=C\_der*(epsilon\_r+r\_iso) for every A in B\_\{r\_iso\}\textasciicircum{}\{icalH\}(0); and (R) set q\_*:=C\_grp*epsilon\_r, set r\_-:=delta\_*-C\_pol*(epsilon\_r*delta\_*+delta\_*\textasciicircum{}2), and set eta\_*:=C\_path*(q\_*+epsilon\_r*q\_*+q\_*\textasciicircum{}2); the guards C\_ch*(epsilon\_r+delta\_*)<=kappa\_ch, C\_pol*(epsilon\_r+delta\_*)<=kappa\_pol, q\_*<r\_-, C\_path*q\_*<=1/4, eta\_*<r\_-, C\_der*(epsilon\_r+r\_iso)<=kappa\_der/4<1, and (1+epsilon\_r)*(1+C\_ch*(epsilon\_r+delta\_*))*r\_iso+q\_*<2delta\_* hold; then, for those same u\_delta,h\_delta,mu,sigma,H,chi\_s,F\_s, there exist a space breve-calU, maps breve-mu:breve-calU x breve-calU->breve-calU and breve-sigma:breve-calU->breve-calU, and maps psi\_s:chi\_s(B\_\{r\_iso\}\textasciicircum{}\{icalH\}(0))->B\_\{r\_iso\}\textasciicircum{}\{icalH\}(0) for s in \{+1,-1\} such that breve-calU=calU\_e/U(1), set breve-e:=[J], breve-mu([U],[V])=[mu(U,V)], breve-sigma([U])=[sigma(U)], (breve-calU,breve-mu,breve-e) is a connected H-space, breve-sigma is a smooth left inversion, sigma(cU)=conj(c)*sigma(U), ||sigma(U)-U\textasciicircum{}dagger||<=C\_grp*epsilon\_r, breve-calU is a connected compact orientable smooth manifold without boundary of real dimension (dim\_C calX)-1, chi\_s:B\_\{r\_iso\}\textasciicircum{}\{icalH\}(0)->calU has chi\_s(0)=sJ and inverse psi\_s=phi\_\{sJ\}\textasciicircum{}par on its image, sigma retains that image, F\_s=psi\_s o sigma o chi\_s and ||D(F\_s-id)(A)+2I||<1, calU intersect B\_\{r\_iso\}(sJ) is contained in chi\_s(B\_\{r\_iso\}\textasciicircum{}\{icalH\}(0)), and ||A-B||<=||chi\_s(A)-chi\_s(B)|| for A,B in B\_\{r\_iso\}\textasciicircum{}\{icalH\}(0); for these same maps, breve-e is an isolated fixed point of breve-sigma, i*reals*J is D-sigma\_J-invariant, ||D-breve-sigma\_\{breve-e\}+I||<1 in the quotient norm, det(I-D-breve-sigma\_\{breve-e\})>0, and the local index of breve-e is +1. \steptag{validated} \\[6pt]
\step{1.1} Choose the universal witnesses e\_quot\textasciicircum{}r and epsilon\_B\textasciicircum{}r furnished by lem-stage1-bound-quotient-left-inversion; under the present stronger hypotheses, establish with the same u\_delta,h\_delta,mu,sigma,H the connected compact orientable smooth quotient H-space, its smooth left inversion, equivariance, and the stated sigma estimate, while the three smooth-upgrade imports identify all smooth structures and maps with the displayed C\textasciicircum{}1 ones. \steptag{validated} \\[6pt]
\step{1.1.1} By (A\_2), ||D\_\{A\textasciicircum{}perp\}f\_V-I||<1, so the Neumann-series criterion makes every such derivative invertible; lem-stage1-smooth-unitary-atlas therefore makes the same g\_V and chi\_V smooth with unchanged points and first derivatives. We now supply the ambient-openness step omitted from the earlier wording. Regard calX as a real finite-dimensional space. The conjugate-linear involution gives the real direct sum calX=calH direct-sum icalH: for every X, X=(X+X\textasciicircum{}dagger)/2+(X-X\textasciicircum{}dagger)/2 with the first summand Hermitian and the second anti-Hermitian, and their intersection is zero. The graph atlas models calU on open subsets of icalH, so dim\_R(calU x B\_\{delta\_*\}\textasciicircum{}\{calH\}(J))=dim\_R(icalH)+dim\_R(calH)=dim\_R(calX). Because (A\_4) says Pi\_delta\_* is a C\textasciicircum{}1 diffeomorphism onto its image as a C\textasciicircum{}1 submanifold of calX, it is a C\textasciicircum{}1 embedding; consequently its ambient differential is injective at every point, and the dimension equality makes that differential a real-linear isomorphism. Fix a source point p and take a graph-product coordinate representative P of Pi\_delta\_* with x\_0 representing p. Put V=DP(x\_0). Continuity of DP permits a ball B\_r(x\_0) in the coordinate domain on which ||V\textasciicircum{}(-1)DP(x)-I||<=1/2. Applying lem-stage1-quantitative-inverse-function, P(B\_r(x\_0)) contains P(x\_0)+V(B\_\{r/2\}(0)), an ambient-open neighborhood of P(x\_0). Since p was arbitrary and S\_delta\_* is the image, S\_delta\_* is open in calX. Now (A\_4), including its C\textasciicircum{}1 inverse, is a bijective C\textasciicircum{}1 local diffeomorphism onto this ambient-open set, so lem-stage1-smooth-polar-inverse makes this same Pi\_delta\_* and the same inverse (u\_delta,h\_delta) smooth. Global definition in (A\_5) implies U bold-dot V and U\textasciicircum{}dagger belong to S\_delta\_*; hence lem-stage1-explicit-smooth-unitary-operations makes the same scalar action, mu, and sigma smooth, with unchanged values/differentials and sigma(cU)=conj(c)*sigma(U). \steptag{validated} \\[4pt]
\step{1.1.2} All assumptions required by lem-stage1-bound-quotient-left-inversion are literally (A\_2),(A\_4),(A\_5),(A\_6) and the first five guards in (R), all present here; take its universal e\_quot\textasciicircum{}r>0 and epsilon\_B\textasciicircum{}r=min\{epsilon\_*\textasciicircum{}r,e\_quot\textasciicircum{}r\}. Its conclusion, for the same u\_delta,h\_delta,mu,sigma,H, gives breve-calU=calU\_e/U(1), breve-e=[J], breve-mu([U],[V])=[mu(U,V)], breve-sigma([U])=[sigma(U)], the connected H-space and smooth-left-inversion assertions, ||sigma(U)-U\textasciicircum{}dagger||<=C\_grp*epsilon\_r, and the connected compact orientable boundaryless smooth-manifold assertion with real dimension dim\_C(calX)-1; the preceding smooth upgrades ensure these are exactly the displayed maps and unchanged C\textasciicircum{}1 differentials. \steptag{validated} \\[4pt]
\step{1.2} For each s in \{+1,-1\}, establish the entire asserted ambient-chart package for chi\_s, psi\_s, and F\_s: smoothness, chi\_s(0)=sJ, psi\_s=phi\_\{sJ\}\textasciicircum{}par on the image, sigma-invariance and the displayed F\_s formula and derivative bound, containment of calU intersect B\_\{r\_iso\}(sJ), and the lower Lipschitz estimate. \steptag{validated} \\[6pt]
\step{1.2.1} For V=sJ, exact unitality and bilinearity give f\_\{sJ\}(0)=0; uniqueness in (A\_2) gives g\_\{sJ\}(0)=0, hence chi\_s(0)=sJ. The graph-coordinate identity phi\_\{sJ\}(chi\_s(A))=A+g\_\{sJ\}(A) shows that the inverse is its anti-Hermitian component psi\_s=phi\_\{sJ\}\textasciicircum{}par. Smoothness follows from lem-stage1-smooth-unitary-atlas. The sigma-invariance and F\_s=psi\_s o sigma o chi\_s are exactly (A\_7), and D(F\_s-id)+2I=DF\_s+I has norm at most C\_der*(epsilon\_r+r\_iso)<=kappa\_der/4<1 by (A\_7) and (R). \steptag{validated} \\[4pt]
\step{1.2.2} If U is in calU with ||U-sJ||<r\_iso, left multiplication by sJ is scalar multiplication by s, so phi:=phi\_\{sJ\}(U)=s(U-sJ). Its anti-Hermitian and Hermitian projections A=(phi-phi\textasciicircum{}dagger)/2 and B=(phi+phi\textasciicircum{}dagger)/2 have norms at most ||phi||<r\_iso<=delta\_*/4<2delta\_*. Since U=sJ bold-dot(J+phi) and U\textasciicircum{}dagger bold-dot U=J, the defining displayed formula gives f\_\{sJ\}(phi)=0. Uniqueness in (A\_2) forces B=g\_\{sJ\}(A), so U=chi\_s(A) with A in B\_\{r\_iso\}\textasciicircum{}\{icalH\}(0); therefore calU intersect B\_\{r\_iso\}(sJ) is contained in the asserted chart image. \steptag{validated} \\[4pt]
\step{1.2.3} For A,B in B\_\{r\_iso\}\textasciicircum{}\{icalH\}(0), exact unitality and bilinearity give chi\_s(A)-chi\_s(B)=s*((A-B)+(g\_\{sJ\}(A)-g\_\{sJ\}(B))). The anti-Hermitian projection P\_parallel(X)=(X-X\textasciicircum{}dagger)/2 is contractive because dagger is an isometry, and it kills the Hermitian g-difference while fixing A-B. Thus ||A-B||=||P\_parallel(chi\_s(A)-chi\_s(B))||<=||chi\_s(A)-chi\_s(B)||. \steptag{validated} \\[4pt]
\step{1.3} At breve-e=[J], prove that i*reals*J is D-sigma\_J-invariant and that the induced quotient differential satisfies ||D-breve-sigma\_\{breve-e\}+I||<1 in the quotient norm. \steptag{validated} \\[6pt]
\step{1.3.1} From sigma(cU)=conj(c)*sigma(U) and sigma(J)=J, differentiating sigma(exp(it)J)=exp(-it)J at t=0 gives D-sigma\_J(iJ)=-iJ, hence i*reals*J is invariant. In the chi\_+ coordinate, the vertical tangent of the U(1)-orbit is the same line i*reals*J (the coordinate of exp(it)J has derivative iJ), so DF\_+(0) acts as -I on that line. \steptag{validated} \\[4pt]
\step{1.3.2} The quotient tangent is T\_\{breve-e\}breve-calU identified with icalH/(i*reals*J), equipped with the quotient norm, and D-breve-sigma\_\{breve-e\} is induced by DF\_+(0). Put R=DF\_+(0)+I. By the preceding vertical calculation R vanishes on i*reals*J, so it induces Rbar=D-breve-sigma\_\{breve-e\}+I. For quotient classes, ||Rbar[x]||\_quot=inf\_v||R(x+v)||<=||R||*inf\_v||x+v||, hence ||Rbar||<=||R||. By (A\_7) and (R), ||R||<=C\_der*(epsilon\_r+r\_iso)<=kappa\_der/4<1, proving the asserted quotient estimate. \steptag{validated} \\[6pt]
\step{1.3.2.1} Correct the quotient-norm step as follows. Let V=i*reals*J and let q:icalH->icalH/V be the quotient map. Since R|\_V=0, the induced map Rbar is well-defined by Rbar(q(x))=q(Rx). Hence, for every x in icalH and every v in V, Rbar(q(x))=q(R(x+v)), and therefore ||Rbar(q(x))||\_quot=||q(R(x+v))||\_quot<=||R(x+v)||<=||R||*||x+v||. Taking the infimum over v in V gives ||Rbar(q(x))||\_quot<=||R||*||q(x)||\_quot, so ||Rbar||<=||R||. Here R=DF\_+(0)+I=D(F\_+-id)(0)+2I, and (A\_7) together with (R) gives ||R||<=C\_der*(epsilon\_r+r\_iso)<=kappa\_der/4<1. Thus ||D-breve-sigma\_\{breve-e\}+I||=||Rbar||<1. No equality between the quotient norm of q(Rx) and the ambient norm of Rx is used. \steptag{validated} \\[4pt]
\step{1.3.2.2} Replace the final norm chain by the following fully non-strict argument. Let V=i*reals*J and q:icalH->icalH/V be the quotient map. With R=DF\_+(0)+I=D(F\_+-id)(0)+2I, the preceding vertical calculation gives R|\_V=0, so Rbar(q(x)):=q(Rx) is well-defined and equals D-breve-sigma\_\{breve-e\}+I. For every x in icalH and v in V, Rbar(q(x))=q(R(x+v)); hence ||Rbar(q(x))||\_quot<=||R(x+v)||<=||R||*||x+v||. Infimizing over v gives ||Rbar(q(x))||\_quot<=||R||*||q(x)||\_quot and therefore ||Rbar||<=||R||. Assumption (A\_7), evaluated at A=0, gives ||R||=||D(F\_+-id)(0)+2I||<=C\_der*(epsilon\_r+r\_iso), while (R) gives C\_der*(epsilon\_r+r\_iso)<=kappa\_der/4<1. Consequently ||D-breve-sigma\_\{breve-e\}+I||=||Rbar||<=||R||<=C\_der*(epsilon\_r+r\_iso)<=kappa\_der/4<1. This uses no unsupported strict inequality and proves the required strict conclusion. \steptag{validated} \\[4pt]
\step{1.4} Prove that breve-e is an isolated fixed point of breve-sigma by applying lem-stage1-quantitative-inverse-function in a local quotient chart to id-breve-sigma. \steptag{validated} \\[6pt]
\step{1.4.1} Take a smooth local chart theta of breve-calU at breve-e with theta(breve-e)=0, write b=theta o breve-sigma o theta\textasciicircum{}\{-1\}, and L=Db(0). The estimate ||L+I||<1 makes T:=I-L=2I-(L+I) invertible by the Neumann series. For q(x):=x-b(x), Dq(0)=T; continuity of Dq permits a ball B\_rho(0) on which ||T\textasciicircum{}\{-1\}Dq(x)-I||<=c<1. Apply lem-stage1-quantitative-inverse-function with V=T: q is injective on this ball. Since sigma(J)=J gives q(0)=0, zero is the only zero of q there, equivalently breve-e is the only fixed point of breve-sigma in the corresponding neighborhood. \steptag{validated} \\[6pt]
\step{1.4.1.1} Let Q=T\_\{breve-e\}breve-calU with its quotient norm and A=D-breve-sigma\_\{breve-e\}. Node 1.3.2 gives ||A+I\_Q||<1. Hence T\_0:=I\_Q-A=2I\_Q-(A+I\_Q)=2(I\_Q-(A+I\_Q)/2) is a Banach-space isomorphism by the Neumann series, since ||(A+I\_Q)/2||<1/2. Now choose any smooth chart theta at breve-e, after shrinking its domain so that breve-sigma maps it into the chart, with theta(breve-e)=0. Put E for its finite-dimensional coordinate space, S=Dtheta\_\{breve-e\}:Q->E, b=theta o breve-sigma o theta\textasciicircum{}\{-1\}, and L=Db(0)=S A S\textasciicircum{}\{-1\}. Then T:=I\_E-L=S T\_0 S\textasciicircum{}\{-1\} is an isomorphism; no bound on ||L+I\_E|| in the coordinate norm is asserted or needed. For q(x)=x-b(x), one has Dq(0)=T. Choose a norm on E (any norm makes finite-dimensional E Banach) and rho>0 with B\_rho(0) in the domain of q. Because x $|-\rangle$ T\textasciicircum{}\{-1\}Dq(x) is continuous and equals I\_E at 0, after decreasing rho there is c<1 such that ||T\textasciicircum{}\{-1\}Dq(x)-I\_E||<=c throughout B\_rho(0). Lem-stage1-quantitative-inverse-function applied with V=T, x\_0=0, and f=q makes q injective on B\_rho(0). Since breve-sigma(breve-e)=breve-e, b(0)=0 and q(0)=0, so q has no other zero there. Thus breve-e is an isolated fixed point of breve-sigma. \steptag{validated} \\[4pt]
\step{1.4.2} Work independently of pending node 1.3.2. Let p:calU\_e->breve-calU be the scalar quotient and V=i*reals*J. In the chi\_+ chart at J, the smooth quotient supplied by lem-stage1-bound-quotient-left-inversion identifies T\_\{breve-e\}breve-calU with Q=icalH/V and its quotient norm; Dp\_J is the quotient map p\_Q:icalH->Q. Since breve-sigma([U])=[sigma(U)], the identity p o sigma=breve-sigma o p and the chain rule show that A:=D-breve-sigma\_\{breve-e\} is induced on Q by B:=DF\_+(0). Indeed F\_+=psi\_+ o sigma o chi\_+ is the local expression of sigma. Scalar equivariance and sigma(J)=J give sigma(exp(it)J)=exp(-it)J; differentiating at zero in the chi\_+ coordinate gives B(iJ)=-iJ, hence B=-I on V. Thus R:=B+I vanishes on V and Rbar(p\_Q x):=p\_Q(Rx) is well-defined, with Rbar=A+I\_Q. For every x in icalH and v in V, Rbar(p\_Q x)=p\_Q(R(x+v)), so ||Rbar(p\_Q x)||\_quot<=||R(x+v)||<=||R||*||x+v||. Infimizing over v yields ||A+I\_Q||=||Rbar||<=||R||. Since R=DF\_+(0)+I=D(F\_+-id)(0)+2I, (A\_7) at zero and (R) give the fully non-strict chain ||R||<=C\_der*(epsilon\_r+r\_iso)<=kappa\_der/4<1. Consequently T\_0:=I\_Q-A=2(I\_Q-(A+I\_Q)/2) is an isomorphism by the Neumann series. Choose a smooth chart theta at breve-e, restricted to the open intersection of its domain with breve-sigma\textasciicircum{}\{-1\} of its domain, and set theta(breve-e)=0, S=Dtheta\_\{breve-e\}, b=theta o breve-sigma o theta\textasciicircum{}\{-1\}, and G(x)=x-b(x). Then DG(0)=I-Db(0)=S T\_0 S\textasciicircum{}\{-1\}=:T is an isomorphism. Equip the finite-dimensional coordinate space with any Banach norm. By continuity of DG, on some ball B\_rho(0) one has ||T\textasciicircum{}\{-1\}DG(x)-I||<=1/2<1. Lem-stage1-quantitative-inverse-function applied to G with V=T makes G injective there. Finally breve-sigma(breve-e)=breve-e because sigma(J)=J, so G(0)=0; injectivity makes 0 its only zero on that ball. Hence breve-e is an isolated fixed point of breve-sigma, without using node 1.3.2. \steptag{validated} \\[4pt]
\step{1.5} Prove det(I-D-breve-sigma\_\{breve-e\})>0 and apply lem-topology-local-index-sign to obtain local index +1. \steptag{validated} \\[6pt]
\step{1.5.1} Let E=D-breve-sigma\_\{breve-e\}+I, so ||E||<1 and I-D-breve-sigma\_\{breve-e\}=2I-E. For 0<=t<=1, 2I-tE is invertible because ||tE/2||<1; its real determinant is continuous and never zero, and at t=0 equals det(2I)>0, hence det(I-D-breve-sigma\_\{breve-e\})>0. The map breve-sigma is smooth on the compact orientable manifold breve-calU, and breve-e is isolated by the preceding node; therefore lem-topology-local-index-sign applies and gives ind(breve-sigma,breve-e)=sgn det(I-D-breve-sigma\_\{breve-e\})=+1. \steptag{validated} \\[4pt]
\end{proofbox}

\end{document}
